\section{Limitations and Open Challenges}
\label{sec:limitations}

\subsection{The decomposition is posited and may not be identifiable}
The five components are motivated by recurring design concerns, not derived as necessary or sufficient constituents of intelligence. Other decompositions may be equally useful. If a vector field is added to one component and subtracted from another, the net dynamics can remain unchanged. Observing trajectories alone consequently need not identify the components. Independent measurements, structural restrictions, or interventions are required before a force label can be given a causal interpretation.

Overlapping cost measures create a related problem. Energy may appear in both computational drag and sustainability, and protection may overlap with risk constraints. A model must explain whether these terms represent distinct timescales, distinct consequences, or deliberate multiple valuation. Without that explanation, apparent advantages may reflect repeated penalties rather than a new mechanism.

\subsection{The mathematical formulation is restricted}
The projected dynamics assumes a continuous state, directly selectable first-order updates, and a fixed convex feasible set. Real systems may have discrete modes, nonconvex constraints, delayed observations, and limited control authority. Existence of an equilibrium, stability, robustness, and tracking do not follow from the force labels. The analytical example establishes those properties only for its scalar linear field and fixed parameters.

A general variational functional has not been derived. It would be misleading to infer a common principle of stationary action from the proposed notation. Extensions to stochastic systems would also need probabilistic definitions of constraint satisfaction and persistence, with explicit disturbance assumptions and time horizons.

\subsection{Empirical and comparative support is absent}
This manuscript reports no new experiment or validated simulation comparison. The companion repository contains illustrative simulations with stipulated objectives and restricted comparators (Section~\ref{subsec:cog_predictions}). H1--H4 describe a prospective evaluation programme; they are not demonstrated effects. Neither the clinical and engineering scenarios nor the biological analogies provide independent validation. The lack of evidence for incremental value over resource-aware alternatives is the principal obstacle to presenting UVIF as an established theory or proven architecture.

The literature discussion is selective. A more comprehensive account should examine viability theory, adaptive and robust control, constrained decision processes, and related dynamical approaches to cognition in depth. Any claim of priority or formal generalization would require that comparison and a precise derivation.

\subsection{Normative and cross-domain limitations}
Admissibility specifications may omit important harms or encode contested priorities. Mathematical feasibility is not a substitute for domain expertise, governance, or ethical judgment. Persistent operation is not inherently desirable, and natural systems provide no automatic standard for acceptable artificial behaviour.

Broad domain translation also risks becoming unfalsifiable if every outcome can be redescribed as a change in force balance. To limit that risk, an application must specify the component measurements, parameter rules, and predicted outcomes before evaluation. A result inconsistent with those commitments should count against that implementation, rather than being absorbed by unconstrained redefinition of its fields.
